\chapter{Loading Pretrained Weights}

Training a large language model from scratch requires enormous computational
resources---hundreds of GPU-hours for even a modest model. Rather than
repeating this expensive process, we can load \emph{pretrained weights}:
parameters that have already been trained on massive text corpora and
published by other researchers. This practice, known as \textbf{transfer
learning}, lets us skip pretraining entirely and immediately use or
fine-tune a model that already understands language.

The most widely used repository for pretrained model weights is
\textbf{HuggingFace}, an open-source platform that hosts thousands of
models. OpenAI released the GPT-2 weights through HuggingFace, making them
accessible to anyone. The general workflow for loading pretrained weights
is:

\begin{enumerate}
  \item Identify the model architecture (e.g., GPT-2 Small with 12 layers,
        768-dimensional embeddings, 12 attention heads).
  \item Download the weight checkpoint from HuggingFace.
  \item Determine the mapping between the checkpoint's parameter names and
        the names in the target architecture.
  \item Handle any differences in how weights are stored (e.g.,
        transpositions, fused matrices).
  \item Assign the weights to the model, skipping any parameters that
        differ in shape or don't have a match.
\end{enumerate}

This chapter walks through each of these steps for GPT-2, explaining the
subtleties that arise when two implementations of the same architecture
store their parameters differently.

\section{Linear Layers and Weight Matrices}

Before diving into weight loading, we need to understand how linear
transformations are represented in neural networks and how their weights
are stored.

A \textbf{linear layer} (also called a fully connected or dense layer)
performs the transformation:
\[
  y = Wx + b
\]
where $x \in \mathbb{R}^{d_{\text{in}}}$ is the input vector,
$W \in \mathbb{R}^{d_{\text{out}} \times d_{\text{in}}}$ is the weight
matrix, and $b \in \mathbb{R}^{d_{\text{out}}}$ is the bias vector.

In PyTorch, \texttt{nn.Linear(in\_features, out\_features)} stores the
weight matrix as a tensor of shape
$[\text{out\_features}, \text{in\_features}]$. This convention comes from
the matrix multiplication $y = Wx$: the weight matrix has one row per
output dimension and one column per input dimension.

\begin{keyconcept}
  The two linear layers in the GPT-2 feed-forward network (FFN) serve
  different roles:
  \begin{itemize}
    \item \textbf{Expansion layer} ($W_1$, called \texttt{ff1}): maps from
          embedding dimension (768) to a hidden dimension (3072). This is a
          $4\times$ expansion that gives the model more capacity to learn
          complex features.
    \item \textbf{Projection layer} ($W_2$, called \texttt{ff2}): maps from
          hidden dimension (3072) back down to embedding dimension (768).
  \end{itemize}
  In GPT-2, the FFN computation is:
  $\text{FFN}(x) = W_2 \cdot \text{GELU}(W_1 x + b_1) + b_2$
\end{keyconcept}

Similarly, the multi-head attention mechanism uses four linear layers in
each transformer block:

\begin{keyconcept}
  Multi-head attention projects the input into three sets of vectors and
  then combines the result:
  \begin{itemize}
    \item $W_Q$ (\texttt{W\_q}): the \textbf{query} projection. Maps input
          from embedding dimension (768) to the total query dimension (768).
          Each attention head uses a 64-dimensional slice.
    \item $W_K$ (\texttt{W\_k}): the \textbf{key} projection. Same
          dimensions as $W_Q$. Keys represent ``what information I hold.''
    \item $W_V$ (\texttt{W\_v}): the \textbf{value} projection. Same
          dimensions. Values represent ``what content I provide when
          attended to.''
    \item $W_O$ (\texttt{W\_o}): the \textbf{output} projection. Maps from
          embedding dimension (768) back to embedding dimension (768),
          combining the outputs of all heads.
  \end{itemize}
  Each of $W_Q$, $W_K$, $W_V$ is stored as an \texttt{nn.Linear(768, 768)}
  with shape $[768, 768]$.
\end{keyconcept}

\section{Conv1D vs.\ Linear: Why the Transpose Matters}

HuggingFace's GPT-2 implementation uses \texttt{Conv1D} layers instead of
\texttt{nn.Linear}. Despite the name, \texttt{Conv1D} here does not
perform convolution---it is simply a linear layer that stores its weight
matrix in a \emph{transposed} layout compared to \texttt{nn.Linear}.

The key difference is in how each layer computes the forward pass:
\begin{itemize}
  \item \texttt{nn.Linear}: $y = x \cdot W^\top + b$, where
        $W$ has shape $[\text{out\_features}, \text{in\_features}]$.
  \item \texttt{Conv1D}: $y = x \cdot W + b$, where
        $W$ has shape $[\text{in\_features}, \text{out\_features}]$.
\end{itemize}

Both compute the same mathematical operation---a linear projection. The
difference is purely in storage layout:

\begin{center}
\begin{tabular}{lll}
\toprule
 & \textbf{HuggingFace (Conv1D)} & \textbf{Standard (nn.Linear)} \\
\midrule
Weight shape & $[\text{in\_features}, \text{out\_features}]$ & $[\text{out\_features}, \text{in\_features}]$ \\
Forward pass & $x \cdot W + b$ & $x \cdot W^\top + b$ \\
Transpose needed when loading? & Yes & No \\
\bottomrule
\end{tabular}
\end{center}

\begin{keyconcept}
  The Conv1D convention comes from early deep learning frameworks where
  a 1D convolution with kernel size 1 is mathematically equivalent to a
  linear layer. HuggingFace inherited this convention from the original
  OpenAI GPT-2 code release. When loading HuggingFace GPT-2 weights into a
  model that uses \texttt{nn.Linear}, every Conv1D weight must be
  transposed.
\end{keyconcept}

To see why the transpose is necessary, consider a concrete example. A
layer projecting from 768 to 3072 dimensions:
\begin{itemize}
  \item HuggingFace stores it as shape $[768, 3072]$.
  \item \texttt{nn.Linear(768, 3072)} expects shape $[3072, 768]$.
\end{itemize}
Loading the HuggingFace weight directly into \texttt{nn.Linear} would cause
a shape mismatch. The fix is to transpose:
\texttt{weight.data = hf\_weight.T}.

\section{Weight Mapping}

The GPT-2 model consists of an embedding layer, 12 transformer blocks, a
final layer norm, and a linear output projection (the LM head). Each
transformer block contains the attention sublayers ($W_Q$, $W_K$, $W_V$,
$W_O$) and the feed-forward sublayers ($W_1$, $W_2$), along with their
biases and layer norm parameters.

HuggingFace's naming convention differs from what a straightforward
implementation might use. The mapping below shows how each parameter in a
standard implementation corresponds to a parameter in the HuggingFace
checkpoint:

\begin{center}
\begin{tabular}{lll}
\toprule
\textbf{Component} & \textbf{Standard name} & \textbf{HuggingFace name} \\
\midrule
Token embedding & \texttt{embedding.token.weight} & \texttt{transformer.wte.weight} \\
Position embedding & \texttt{embedding.positional.weight} & \texttt{transformer.wpe.weight} \\
Query projection & \texttt{blocks[$i$].attn.W\_q.weight} & \texttt{h[$i$].attn.c\_attn.weight}[:, :768] \\
Key projection & \texttt{blocks[$i$].attn.W\_k.weight} & \texttt{h[$i$].attn.c\_attn.weight}[:, 768:1536] \\
Value projection & \texttt{blocks[$i$].attn.W\_v.weight} & \texttt{h[$i$].attn.c\_attn.weight}[:, 1536:] \\
Output projection & \texttt{blocks[$i$].attn.W\_o.weight} & \texttt{h[$i$].attn.c\_proj.weight} \\
FFN expansion & \texttt{blocks[$i$].ff.ff1.weight} & \texttt{h[$i$].mlp.c\_fc.weight} \\
FFN projection & \texttt{blocks[$i$].ff.ff2.weight} & \texttt{h[$i$].mlp.c\_proj.weight} \\
\bottomrule
\end{tabular}
\end{center}

There are two important differences beyond naming:

\subsection{Fused QKV Matrix}

HuggingFace stores the query, key, and value projections as a single fused
matrix:

\begin{keyconcept}
  HuggingFace combines $W_Q$, $W_K$, and $W_V$ into one
  \texttt{c\_attn} matrix of shape $[768, 2304]$. This is the
  concatenation of $W_Q$, $W_K$, and $W_V$ along the output dimension:
  \[
    \texttt{c\_attn.weight} = [W_Q \;||\; W_K \;||\; W_V]
  \]
  where each sub-matrix has shape $[768, 768]$ and $3 \times 768 = 2304$.
  To extract the individual projections, we split along the second axis.
\end{keyconcept}

\subsection{Transpositions}

Every weight listed in the table above comes from a \texttt{Conv1D} layer
in HuggingFace's implementation, which means every one of these weights
must be transposed when loading into an \texttt{nn.Linear}-based model.

\section{Loading the Weights}

Loading pretrained weights requires iterating through the checkpoint's
parameter names, matching each to the correct parameter in the target
model, and applying any necessary transformations (transpositions,
splits).

\begin{pseudocode}[Load HuggingFace GPT-2 Weights]
def load_gpt2_weights(model: GPTModel, hf_params: dict) -> None:
    # hf_params: flat dict mapping parameter names to tensors
    # Number of layers extracted from the checkpoint
    n_layers = hf_params["transformer.h.__0__"].shape  # or infer from keys
    for name, param in model.named_parameters():
        # Skip biases and layer norms for brevity
        # Match name patterns to HuggingFace equivalents
        if name == "embedding.token.weight":
            param.data = hf_params["transformer.wte.weight"]
        elif name == "embedding.positional.weight":
            param.data = hf_params["transformer.wpe.weight"]
        elif "attn.W_q" in name:
            # Slice columns 0:768 from the fused c_attn, then transpose
            # Shape: [768, 768] (after transpose) -> [768, 768]
            W_q = hf_params[hf_name][:, :768].T
            param.data = W_q
        elif "attn.W_k" in name:
            # Slice columns 768:1536 from the fused c_attn, then transpose
            W_k = hf_params[hf_name][:, 768:1536].T
            param.data = W_k
        elif "attn.W_v" in name:
            # Slice columns 1536: from the fused c_attn, then transpose
            W_v = hf_params[hf_name][:, 1536:].T
            param.data = W_v
        elif "attn.W_o" in name:
            # c_proj has shape [768, 768], transpose for nn.Linear
            param.data = hf_params[hf_name].T
        elif "ff.ff1" in name:
            # c_fc has shape [768, 3072], transpose for nn.Linear
            # After transpose: [3072, 768] matching nn.Linear(768, 3072)
            param.data = hf_params[hf_name].T
        elif "ff.ff2" in name:
            # c_proj has shape [3072, 768], transpose for nn.Linear
            # After transpose: [768, 3072] matching nn.Linear(3072, 768)
            param.data = hf_params[hf_name].T
\end{pseudocode}

\begin{keyconcept}
  The weight loading process involves three types of transformations:
  \begin{itemize}
    \item \textbf{Direct copy}: token and positional embeddings transfer
          directly with no transformation.
    \item \textbf{Transpose}: all linear and Conv1D weights require
          transposition because of the storage layout difference.
    \item \textbf{Split + transpose}: the fused QKV matrix must first be
          split into three slices along the output dimension, then each
          slice is transposed.
  \end{itemize}
\end{keyconcept}

A robust implementation should also handle parameters that don't have a
match---the final layer norm, biases, and any extra parameters should be
skipped with a warning rather than causing an error.

\section{Weight Tying}

In GPT-2, the token embedding matrix and the output projection (the
\emph{language model head}, or LM head) share the same weight tensor:

\[
  \texttt{lm\_head.weight} = \texttt{embedding.token.weight}
\]

This practice, called \textbf{weight tying} (or weight sharing), means the
same matrix is used for two different purposes:

\begin{enumerate}
  \item \textbf{At input}: mapping token IDs to embedding vectors (looking
        up rows of the matrix).
  \item \textbf{At output}: projecting the final hidden state to vocabulary
        logits (matrix multiplication with the hidden state).
\end{enumerate}

\begin{keyconcept}
  Weight tying saves approximately 38M parameters for GPT-2 Small
  (vocabulary size 50,257 $\times$ embedding dimension 768 $= 38,597,376$
  parameters, roughly one-third of the model's 124M total). Beyond saving
  memory, weight tying has a theoretical motivation: the same
  representations used to encode tokens should be useful for decoding them.
  If a token is embedded near certain vectors, the output projection should
  assign high probability to that same token when those vectors appear.
  This symmetry has been shown to improve model quality, especially in
  early training.
\end{keyconcept}

When loading HuggingFace weights, the LM head matrix does not appear as a
separate entry in the checkpoint---it is simply the token embedding matrix
reused. Any implementation must explicitly tie these weights:

\begin{pseudocode}[Weight Tying]
def tie_weights(model: GPTModel) -> None:
    # lm_head is a nn.Linear layer projecting from 768 to 50257
    # embedding.token is an nn.Embedding of shape [50257, 768]
    model.lm_head.weight = model.embedding.token.weight
    # Now both parameters point to the same underlying tensor
\end{pseudocode}

\section{Validating Loaded Weights}

After loading weights, it is crucial to verify that the model works
correctly. A model with improperly loaded weights---for example, if a
transposition was forgotten or if Q and V were swapped---will produce
gibberish rather than coherent text.

The simplest validation is to generate text from a known prompt and check
that the output is coherent:

\begin{pseudocode}[Validate Loaded Weights via Text Generation]
def validate_weights(model: GPTModel, tokenizer: Tokenizer) -> None:
    # The model should produce coherent English text
    prompt = "The capital of France is"
    token_ids = tokenizer.encode(prompt)
    # [batch, seq_len] = [1, 5]
    input_tensor = torch.tensor([token_ids])

    # Generate 20 new tokens using greedy decoding
    for _ in range(20):
        # logits shape: [1, seq_len, vocab_size]
        logits = model(input_tensor)
        # Take logits at the last position only
        next_token_id = logits[:, -1, :].argmax(dim=-1)  # [1]
        input_tensor = torch.cat([input_tensor, next_token_id.unsqueeze(1)], dim=1)

    output_text = tokenizer.decode(input_tensor[0])
    # Expected: coherent continuation like " Paris, which is the..."
    print(output_text)
\end{pseudocode}

\begin{keyconcept}
  \textbf{Common bugs when loading weights:}
  \begin{itemize}
    \item \textbf{Forgetting to transpose}: the model runs but produces
          gibberish because every linear layer computes $x \cdot W$ instead
          of $x \cdot W^\top$.
    \item \textbf{Swapping Q, K, V}: splitting the fused matrix incorrectly
          results in nonsensical attention patterns.
    \item \textbf{Missing weight tying}: the LM head uses random weights
          while the embedding uses pretrained weights, so the model cannot
          predict tokens correctly.
    \item \textbf{Shape mismatches}: silently assigning a wrong-shaped tensor
          (e.g., not splitting the fused QKV) causes dimension errors.
  \end{itemize}
\end{keyconcept}

Another useful check is to verify that the loaded parameters match the
expected shapes and have reasonable statistics (mean near zero, standard
deviation near 0.02 for GPT-2):

\begin{pseudocode}[Verify Weight Shapes and Statistics]
def verify_weights(model: GPTModel) -> None:
    for name, param in model.named_parameters():
        # Check shape matches expected dimensions
        print(f"{name}: shape={param.shape}")
        # GPT-2 weights initialized from N(0, 0.02)
        mean = param.data.mean().item()
        std = param.data.std().item()
        # Embedding weights have different scale; skip them
        if "embedding" not in name:
            assert abs(mean) < 0.01, f"{name} has suspicious mean: {mean}"
            # std should be close to 0.02 for trained weights
            assert 0.005 < std < 0.1, f"{name} has suspicious std: {std}"
\end{pseudocode}

\section{Perplexity Evaluation}

Once the pretrained weights are loaded and validated, we need a way to
quantitatively measure how well the model predicts text. The standard
metric for evaluating language models is \textbf{perplexity}.

\subsection{What Is Perplexity?}

Perplexity measures how ``surprised'' the model is by the text it sees. A
model that always assigns high probability to the correct next token has
low perplexity; a model that is frequently surprised has high perplexity.

Formally, given a sequence of tokens $w_1, w_2, \ldots, w_N$, the model
assigns each token a probability conditioned on the preceding tokens:
$P(w_i \mid w_1, \ldots, w_{i-1})$. Perplexity is defined as:

\[
  \text{Perplexity} = \exp\!\left(\frac{1}{N}\sum_{i=1}^{N} -\log P(w_i \mid w_{<i})\right)
  = \exp(\text{average cross-entropy loss per token})
\]

\begin{keyconcept}
  Perplexity can be interpreted as ``on average, how many tokens is the
  model equally uncertain about at each step.'' A perplexity of 30
  means the model is as uncertain as if it were choosing uniformly among
  30 tokens at each position.

  \textbf{Lower perplexity is better.} A perfect model that always assigns
  probability 1 to the correct next token would have perplexity 1. GPT-2
  Small achieves a perplexity around 30--40 on standard benchmarks like
  WikiText-103.
\end{keyconcept}

\subsection{Computing Perplexity}

Computing perplexity from a trained model involves three steps: tokenize
the evaluation text, run the model on the token sequences, and average the
cross-entropy loss across all tokens.

\begin{pseudocode}[Compute Perplexity]
def compute_perplexity(model: GPTModel, texts: list[str],
                       tokenizer: Tokenizer,
                       max_length: int = 1024) -> float:
    total_loss = 0.0
    total_tokens = 0
    for text in texts:
        # Tokenize and truncate to max_length
        token_ids = tokenizer.encode(text)[:max_length]
        # input: [seq_len] = [max_length]
        input_tensor = torch.tensor([token_ids])
        # targets: shifted by 1 position
        targets = input_tensor[:, 1:]   # [1, seq_len - 1]
        logits = model(input_tensor)[:, :-1, :]  # [1, seq_len - 1, vocab_size]

        # Cross-entropy loss averaged over tokens
        loss = cross_entropy(logits.reshape(-1, vocab_size),
                             targets.reshape(-1))
        total_loss += loss.item() * (len(token_ids) - 1)
        total_tokens += len(token_ids) - 1

    avg_loss = total_loss / total_tokens
    perplexity = math.exp(avg_loss)
    return perplexity
\end{pseudocode}

\begin{keyconcept}
  \textbf{Context length matters}: longer contexts give lower perplexity
  because the model has more information to condition on. Always compare
  perplexity at the same context length. A model evaluated with 1024-token
  contexts will appear better than the same model evaluated with 128-token
  contexts, even though the model itself is identical.
\end{keyconcept}

When computing perplexity, it is important to handle long texts properly.
The standard approach is to slide a window of length \texttt{max\_length}
across the text with some stride, computing loss only on the tokens that
fit within each window. This ensures every token is evaluated with the
maximum possible context.

\section*{Review Questions}

\begin{reviewquestion}[1]
  HuggingFace GPT-2 uses Conv1D with weight shape
  $[\text{in}, \text{out}]$ while PyTorch's \texttt{nn.Linear} uses
  $[\text{out}, \text{in}]$. What operation must you perform when
  loading weights?
\end{reviewquestion}

\begin{reviewquestion}[2]
  Explain weight tying between the token embedding and the LM head.
  Why does it make sense for these two matrices to share parameters?
\end{reviewquestion}

\begin{reviewquestion}[3]
  You evaluate the same model at context length 128 and context length
  1024. Which evaluation gives lower perplexity, and why?
\end{reviewquestion}

\begin{reviewquestion}[4]
  HuggingFace stores $W_Q$, $W_K$, and $W_V$ as a single fused matrix
  of shape $[768, 2304]$. How would you extract the individual weight
  matrices, and what is the shape of each after extraction?
\end{reviewquestion}

\begin{reviewquestion}[5]
  Suppose you forgot to transpose a Conv1D weight when loading it into an
  \texttt{nn.Linear} layer. Will the model crash, produce garbage, or
  produce slightly degraded output? Explain your reasoning.
\end{reviewquestion}

\begin{reviewquestion}[6]
  How many parameters does weight tying save for GPT-2 Small? Express
  your answer both as a count and as a fraction of total model parameters.
\end{reviewquestion}

\begin{reviewquestion}[7]
  Why does the GPT-2 FFN expand the hidden dimension from 768 to 3072
  before projecting back down? What would happen if the hidden dimension
  were equal to the embedding dimension (i.e., no expansion)?
\end{reviewquestion}

\begin{reviewquestion}[8]
  A model achieves a perplexity of 1.0 on a dataset. What does this mean
  about the model's predictions? Is this realistic in practice?
\end{reviewquestion}

\begin{reviewquestion}[9]
  You load GPT-2 weights but forget to tie the LM head weights to the
  embedding weights. The LM head is now using its randomly initialized
  weights. What will the model's output look like, and why?
\end{reviewquestion}

\begin{reviewquestion}[10]
  When computing perplexity on a long document, why is it insufficient to
  simply truncate the text to \texttt{max\_length}? Describe a better
  approach.
\end{reviewquestion}

\begin{reviewquestion}[11]
  The query, key, and value projections all map from 768 to 768
  dimensions. Could they share a single weight matrix? What would be the
  consequence?
\end{reviewquestion}

\begin{reviewquestion}[12]
  Why is Conv1D called ``Conv1D'' when it performs a linear (matrix
  multiplication) operation? What is the historical reason?
\end{reviewquestion}

\begin{reviewquestion}[13]
  Explain the mathematical relationship between cross-entropy loss and
  perplexity. If a model's average cross-entropy loss is 3.5, what is its
  perplexity?
\end{reviewquestion}

\section*{Problems}

\begin{problem}[1]
  Write a function \texttt{load\_gpt2\_weights(model, checkpoint\_path)} that
  loads GPT-2 weights from a HuggingFace checkpoint file
  (\texttt{.bin} or \texttt{.safetensors}). Your function should:
  (a) iterate through all parameters in the checkpoint,
  (b) match each to the correct parameter in the model,
  (c) apply the necessary transpositions for Conv1D weights, and
  (d) handle the fused QKV matrix by splitting it into three separate
  weight matrices. Print a summary of loaded vs.\ skipped parameters.
\end{problem}

\begin{problem}[2]
  Verify weight tying in a loaded GPT-2 model. Write code that:
  (a) loads the pretrained weights including weight tying,
  (b) checks that \texttt{model.lm\_head.weight} and
  \texttt{model.embedding.token.weight} point to the same tensor in memory
  (use \texttt{id()} or \texttt{torch.equal()}),
  (c) modifies one entry in the token embedding and verifies that the
  corresponding entry in the LM head also changes.
\end{problem}

\begin{problem}[3]
  Implement a perplexity evaluation function that processes an arbitrary
  text dataset. Use a sliding window approach with configurable stride.
  Compare the perplexity of a model with correctly loaded GPT-2 weights
  against the same model with randomly initialized weights. How many
  orders of magnitude different are the perplexities?
\end{problem}

\begin{problem}[4]
  Write a diagnostic function \texttt{diagnose\_weight\_loading(model)} that
  checks each layer of a loaded GPT-2 model for common loading errors:
  (a) verify that weight matrices have the expected shapes,
  (b) check that means are near zero and standard deviations are near 0.02,
  (c) detect any weights that are still at their random initialization
  (by checking if the standard deviation matches the initialization
  distribution rather than the pretrained distribution), and
  (d) generate a short text sample and print the first few predicted tokens
  to verify coherence.
\end{problem}

\begin{problem}[5]
  The fused QKV matrix in HuggingFace GPT-2 has shape $[768, 2304]$. Write
  a function that given this matrix, splits it into $W_Q$, $W_K$, and
  $W_V$, transposes each for use in \texttt{nn.Linear}, and assigns them
  to the corresponding model parameters. Then verify correctness by:
  (a) reconstructing the fused matrix from the individual weights and
  confirming it matches the original, and (b) checking that a forward pass
  through the attention layer produces the same output whether you use
  the fused or split weights.
\end{problem}